
\documentclass{article}
\usepackage{colortbl}
\usepackage{makecell}
\usepackage{multirow}
\usepackage{supertabular}

\begin{document}

\newcounter{utterance}

\centering \large Interaction Transcript for game `clean\_up', experiment `0\_easy\_7obj\_it', episode 0 with glm{-}5.2.
\vspace{24pt}

{ \footnotesize  \setcounter{utterance}{1}
\setlength{\tabcolsep}{0pt}
\begin{supertabular}{c@{$\;$}|p{.15\linewidth}@{}p{.15\linewidth}p{.15\linewidth}p{.15\linewidth}p{.15\linewidth}p{.15\linewidth}}
    \# & $\;$A & \multicolumn{4}{c}{Game Master} & $\;\:$B\\
    \hline

    \theutterance \stepcounter{utterance}  
    & & \multicolumn{4}{p{0.6\linewidth}}{
        \cellcolor[rgb]{0.9,0.9,0.9}{
            \makecell[{{p{\linewidth}}}]{
                \texttt{\tiny{[P1$\langle$GM]}}
                \texttt{Sono il tuo maestro di gioco e stai giocando un gioco collaborativo con la seguente griglia come tabellone di gioco:} \\
\\ 
\texttt{```} \\
\texttt{ 1234567} \\
\texttt{╔═╤═══╤═╗} \\
\texttt{║◌│◌DT│◌║ 1} \\
\texttt{╟─┼──┬┘◌║ 2} \\
\texttt{║◌│◌◌│◌◌║ 3} \\
\texttt{╟─┴──┘◌O║ 4} \\
\texttt{║M◌◌P◌◌◌║ 5} \\
\texttt{║◌◌◌◌◌◌◌║ 6} \\
\texttt{║S◌◌◌◌◌A║ 7} \\
\texttt{╚═══════╝} \\
\\ 
\texttt{```} \\
\\ 
\texttt{* Il bordo superiore mostra le coordinate x che aumentano verso destra, e il bordo destro le coordinate y che aumentano verso il basso.} \\
\texttt{* I seguenti oggetti sono posizionati casualmente sulla tua griglia: 'P', 'O', 'T', 'S', 'D', 'A', 'M'.} \\
\\ 
\\ 
\texttt{L'altro giocatore vede una variante del tabellone di gioco, dove gli oggetti sono posizionati in diverse posizioni casuali. Non puoi vedere il tabellone dell'altro giocatore e lui non può vedere il tuo.} \\
\\ 
\\ 
\texttt{**Obiettivo:**} \\
\\ 
\texttt{Entrambi i giocatori devono spostare gli oggetti sui rispettivi sfondi in modo che oggetti identici finiscano alle stesse coordinate. Devi comunicare con l'altro giocatore per concordare una configurazione obiettivo comune.} \\
\\ 
\\ 
\texttt{**Regole:**} \\
\\ 
\texttt{* In ogni turno, puoi inviare esattamente uno dei seguenti due comandi:} \\
\texttt{ 1. `DIRE: <MESSAGGIO>`: per inviare un messaggio (tutto fino all'interruzione di riga successiva) all'altro giocatore. Lo inoltrerò al tuo partner.} \\
\texttt{2. `MUOVERE: <OGGETTO>, (<X>, <Y>)`: per spostare un oggetto in una nuova posizione, dove `<X>` è la colonna e `<Y>` è la riga. Ti informerò se il tuo spostamento ha avuto successo o meno.} \\
\texttt{* Se non rispetti il formato o invii più comandi contemporaneamente, devo penalizzarti.} \\
\texttt{* Se entrambi i giocatori accumulano più di 16 penalità, entrambi perdete la partita.} \\
\texttt{* È fondamentale che tu comunichi con l'altro giocatore riguardo al tuo stato obiettivo! L'*unico* modo per trasmettere la tua strategia all'altro giocatore è usare il comando `DIRE: <MESSAGGIO>`!} \\
\\ 
\\ 
\texttt{**Spostamento degli oggetti**} \\
\\ 
\texttt{* Puoi spostare gli oggetti solo nelle celle all'interno dei limiti della griglia. La cella di destinazione deve essere vuota, cioè deve contenere solo il simbolo '◌'.} \\
\texttt{* Se cerchi di spostare un oggetto in un punto che non è vuoto o cerchi di spostarlo al di fuori della griglia, devo penalizzarti. Avrai un altro tentativo.} \\
\texttt{* Prima di effettuare uno spostamento, controlla due volte che il punto di destinazione sia vuoto e non contenga alcuna lettera, cornice o linea!} \\
\\ 
\\ 
\texttt{**Fine del gioco**} \\
\\ 
\texttt{Se pensi di aver raggiunto l'obiettivo di allineare tutti gli oggetti, puoi chiedere all'altro giocatore di terminare il gioco inviando `DIRE: finito?`. Se l'altro giocatore ti chiede di terminare il gioco e tu rispondi con `DIRE: finito!`, il gioco terminerà.} \\
\\ 
\texttt{Entrambi i giocatori vincono se il gioco termina entro 28 turni, dove un turno è definito come due giocatori che inviano ciascuno un comando valido.} \\
\\ 
\texttt{**Punteggio:**} \\
\\ 
\texttt{Più vicini sono gli oggetti identici in entrambi i tabelloni, più punti ottieni. Le penalità riducono i tuoi punti. Riuscirai a battere il record?} \\
\\ 
\\ 
\texttt{Invia un messaggio all'altro giocatore per iniziare la partita!} \\
            }
        }
    }
    & \\ \\

    \theutterance \stepcounter{utterance}  
    & \multicolumn{4}{p{0.6\linewidth}}{
        \cellcolor[rgb]{0.9,0.9,0.9}{
            \makecell[{{p{\linewidth}}}]{
                \texttt{\tiny{[P1$\rangle$GM]}}
                \texttt{DIRE: Ciao! Ho letto la mia griglia. Propongo come configurazione obiettivo finale queste coordinate per ogni oggetto:} \\
\texttt{P=(4,5), S=(1,7), D=(3,1), T=(4,1), A=(7,7), O=(7,4), M=(1,5).} \\
\texttt{Sono tutte celle vuote nel mio tabellone. Dimmi se per te va bene o proponi altre posizioni!} \\
            }
        }
    }
    & & \\ \\

    \theutterance \stepcounter{utterance}  
    & & & \multicolumn{2}{p{0.3\linewidth}}{
        \cellcolor[rgb]{0.9,0.9,0.9}{
            \makecell[{{p{\linewidth}}}]{
                \texttt{\tiny{[GM$|$GM]}}
                \texttt{DIRE: Ciao! Ho letto la mia griglia. Propongo come configurazione obiettivo finale queste coordinate per ogni oggetto:} \\
\texttt{P=(4,5), S=(1,7), D=(3,1), T=(4,1), A=(7,7), O=(7,4), M=(1,5).} \\
\texttt{Sono tutte celle vuote nel mio tabellone. Dimmi se per te va bene o proponi altre posizioni!} \\
            }
        }
    }
    & & \\ \\

    \theutterance \stepcounter{utterance}  
    & & & \multicolumn{2}{p{0.3\linewidth}}{
        \cellcolor[rgb]{0.9,0.9,0.9}{
            \makecell[{{p{\linewidth}}}]{
                \texttt{\tiny{[GM$|$GM]}}
                \texttt{Invalid format: tail is not empty: '} \\
\texttt{P=(4,5), S=(1,7), D=(3,1), T=(4,1), A=(7,7), O=(7,4), M=(1,5).} \\
\texttt{Sono tutte celle vuote nel mio tabellone. Dimmi se per te va bene o proponi altre posizioni!} \\
            }
        }
    }
    & & \\ \\

    \theutterance \stepcounter{utterance}  
    & & \multicolumn{4}{p{0.6\linewidth}}{
        \cellcolor[rgb]{0.9,0.9,0.9}{
            \makecell[{{p{\linewidth}}}]{
                \texttt{\tiny{[P1$\langle$GM]}}
                \texttt{Penalità: Il tuo messaggio non deve contenere nulla dopo il comando!} \\
\\ 
\texttt{Assicurati che la tua risposta contenga solo `DIRE: <MESSAGGIO>` o `MUOVERE: <OGGETTO>, (<X>, <Y>)`, e nient'altro!} \\
\\ 
\texttt{Avete accumulato collettivamente 1 di 16 penalità.} \\
\texttt{Per favore riprova!} \\
            }
        }
    }
    & \\ \\

    \theutterance \stepcounter{utterance}  
    & \multicolumn{4}{p{0.6\linewidth}}{
        \cellcolor[rgb]{0.9,0.9,0.9}{
            \makecell[{{p{\linewidth}}}]{
                \texttt{\tiny{[P1$\rangle$GM]}}
                \texttt{DIRE: Ciao! Propongo queste coordinate finali per gli oggetti: P=(4,5), S=(1,7), D=(3,1), T=(4,1), A=(7,7), O=(7,4), M=(1,5). Dimmi se per te va bene!} \\
            }
        }
    }
    & & \\ \\

    \theutterance \stepcounter{utterance}  
    & & & \multicolumn{2}{p{0.3\linewidth}}{
        \cellcolor[rgb]{0.9,0.9,0.9}{
            \makecell[{{p{\linewidth}}}]{
                \texttt{\tiny{[GM$|$GM]}}
                \texttt{DIRE: Ciao! Propongo queste coordinate finali per gli oggetti: P=(4,5), S=(1,7), D=(3,1), T=(4,1), A=(7,7), O=(7,4), M=(1,5). Dimmi se per te va bene!} \\
            }
        }
    }
    & & \\ \\

    \theutterance \stepcounter{utterance}  
    & & \multicolumn{4}{p{0.6\linewidth}}{
        \cellcolor[rgb]{0.9,0.9,0.9}{
            \makecell[{{p{\linewidth}}}]{
                \texttt{\tiny{[GM$\rangle$P2]}}
                \texttt{Sono il tuo maestro di gioco e stai giocando un gioco collaborativo con la seguente griglia come tabellone di gioco:} \\
\\ 
\texttt{```} \\
\texttt{ 1234567} \\
\texttt{╔═╤═══╤═╗} \\
\texttt{║◌│◌T◌│◌║ 1} \\
\texttt{╟─┼──┬┘◌║ 2} \\
\texttt{║A│M◌│◌◌║ 3} \\
\texttt{╟─┴──┘◌S║ 4} \\
\texttt{║◌◌◌◌◌◌◌║ 5} \\
\texttt{║◌P◌◌◌◌◌║ 6} \\
\texttt{║◌◌O◌D◌◌║ 7} \\
\texttt{╚═══════╝} \\
\\ 
\texttt{```} \\
\\ 
\texttt{* Il bordo superiore mostra le coordinate x che aumentano verso destra, e il bordo destro le coordinate y che aumentano verso il basso.} \\
\texttt{* I seguenti oggetti sono posizionati casualmente sulla tua griglia: 'P', 'O', 'T', 'S', 'D', 'A', 'M'.} \\
\\ 
\\ 
\texttt{L'altro giocatore vede una variante del tabellone di gioco, dove gli oggetti sono posizionati in diverse posizioni casuali. Non puoi vedere il tabellone dell'altro giocatore e lui non può vedere il tuo.} \\
\\ 
\\ 
\texttt{**Obiettivo:**} \\
\\ 
\texttt{Entrambi i giocatori devono spostare gli oggetti sui rispettivi sfondi in modo che oggetti identici finiscano alle stesse coordinate. Devi comunicare con l'altro giocatore per concordare una configurazione obiettivo comune.} \\
\\ 
\\ 
\texttt{**Regole:**} \\
\\ 
\texttt{* In ogni turno, puoi inviare esattamente uno dei seguenti due comandi:} \\
\texttt{ 1. `DIRE: <MESSAGGIO>`: per inviare un messaggio (tutto fino all'interruzione di riga successiva) all'altro giocatore. Lo inoltrerò al tuo partner.} \\
\texttt{2. `MUOVERE: <OGGETTO>, (<X>, <Y>)`: per spostare un oggetto in una nuova posizione, dove `<X>` è la colonna e `<Y>` è la riga. Ti informerò se il tuo spostamento ha avuto successo o meno.} \\
\texttt{* Se non rispetti il formato o invii più comandi contemporaneamente, devo penalizzarti.} \\
\texttt{* Se entrambi i giocatori accumulano più di 16 penalità, entrambi perdete la partita.} \\
\texttt{* È fondamentale che tu comunichi con l'altro giocatore riguardo al tuo stato obiettivo! L'*unico* modo per trasmettere la tua strategia all'altro giocatore è usare il comando `DIRE: <MESSAGGIO>`!} \\
\\ 
\\ 
\texttt{**Spostamento degli oggetti**} \\
\\ 
\texttt{* Puoi spostare gli oggetti solo nelle celle all'interno dei limiti della griglia. La cella di destinazione deve essere vuota, cioè deve contenere solo il simbolo '◌'.} \\
\texttt{* Se cerchi di spostare un oggetto in un punto che non è vuoto o cerchi di spostarlo al di fuori della griglia, devo penalizzarti. Avrai un altro tentativo.} \\
\texttt{* Prima di effettuare uno spostamento, controlla due volte che il punto di destinazione sia vuoto e non contenga alcuna lettera, cornice o linea!} \\
\\ 
\\ 
\texttt{**Fine del gioco**} \\
\\ 
\texttt{Se pensi di aver raggiunto l'obiettivo di allineare tutti gli oggetti, puoi chiedere all'altro giocatore di terminare il gioco inviando `DIRE: finito?`. Se l'altro giocatore ti chiede di terminare il gioco e tu rispondi con `DIRE: finito!`, il gioco terminerà.} \\
\\ 
\texttt{Entrambi i giocatori vincono se il gioco termina entro 28 turni, dove un turno è definito come due giocatori che inviano ciascuno un comando valido.} \\
\\ 
\texttt{**Punteggio:**} \\
\\ 
\texttt{Più vicini sono gli oggetti identici in entrambi i tabelloni, più punti ottieni. Le penalità riducono i tuoi punti. Riuscirai a battere il record?} \\
\\ 
\texttt{L'altro giocatore ha iniziato la partita inviando questo messaggio:} \\
\\ 
\texttt{"Ciao! Propongo queste coordinate finali per gli oggetti: P=(4,5), S=(1,7), D=(3,1), T=(4,1), A=(7,7), O=(7,4), M=(1,5). Dimmi se per te va bene!"} \\
\\ 
\texttt{Qual è il tuo primo comando?} \\
            }
        }
    }
    & \\ \\

    \theutterance \stepcounter{utterance}  
    & & & \multicolumn{4}{p{0.6\linewidth}}{
        \cellcolor[rgb]{0.9,0.9,0.9}{
            \makecell[{{p{\linewidth}}}]{
                \texttt{\tiny{[GM$\langle$P2]}}
                \texttt{DIRE: La tua proposta va benissimo! Confermo le coordinate finali: P=(4,5), S=(1,7), D=(3,1), T=(4,1), A=(7,7), O=(7,4), M=(1,5). Inizio a spostare M da (2,3) a (1,5).} \\
            }
        }
    }
     \\ \\

    \theutterance \stepcounter{utterance}  
    & & & \multicolumn{2}{p{0.3\linewidth}}{
        \cellcolor[rgb]{0.9,0.9,0.9}{
            \makecell[{{p{\linewidth}}}]{
                \texttt{\tiny{[GM$|$GM]}}
                \texttt{DIRE: La tua proposta va benissimo! Confermo le coordinate finali: P=(4,5), S=(1,7), D=(3,1), T=(4,1), A=(7,7), O=(7,4), M=(1,5). Inizio a spostare M da (2,3) a (1,5).} \\
            }
        }
    }
    & & \\ \\

    \theutterance \stepcounter{utterance}  
    & & \multicolumn{4}{p{0.6\linewidth}}{
        \cellcolor[rgb]{0.9,0.9,0.9}{
            \makecell[{{p{\linewidth}}}]{
                \texttt{\tiny{[P1$\langle$GM]}}
                \texttt{Il tuo messaggio è stato trasmesso all'altro giocatore.} \\
\\ 
\\ 
\texttt{Stai giocando il round 1 di un massimo di 28.} \\
\texttt{Avete accumulato collettivamente 1 di 16 penalità.} \\
\texttt{L'altro giocatore ha inviato questo messaggio:} \\
\\ 
\\ 
\texttt{"La tua proposta va benissimo! Confermo le coordinate finali: P=(4,5), S=(1,7), D=(3,1), T=(4,1), A=(7,7), O=(7,4), M=(1,5). Inizio a spostare M da (2,3) a (1,5)."} \\
\\ 
\\ 
\texttt{Qual è il tuo prossimo comando?} \\
            }
        }
    }
    & \\ \\

    \theutterance \stepcounter{utterance}  
    & \multicolumn{4}{p{0.6\linewidth}}{
        \cellcolor[rgb]{0.9,0.9,0.9}{
            \makecell[{{p{\linewidth}}}]{
                \texttt{\tiny{[P1$\rangle$GM]}}
                \texttt{MUOVERE: T, (4,1)} \\
            }
        }
    }
    & & \\ \\

    \theutterance \stepcounter{utterance}  
    & & & \multicolumn{2}{p{0.3\linewidth}}{
        \cellcolor[rgb]{0.9,0.9,0.9}{
            \makecell[{{p{\linewidth}}}]{
                \texttt{\tiny{[GM$|$GM]}}
                \texttt{MUOVERE: T, (4,1)} \\
            }
        }
    }
    & & \\ \\

    \theutterance \stepcounter{utterance}  
    & & & \multicolumn{2}{p{0.3\linewidth}}{
        \cellcolor[rgb]{0.9,0.9,0.9}{
            \makecell[{{p{\linewidth}}}]{
                \texttt{\tiny{[GM$|$GM]}}
                \texttt{Penalità: (4,1) non è vuoto, ma contiene 'D'.} \\
\\ 
\texttt{Avete accumulato collettivamente 2 di 16 penalità.} \\
\texttt{Per favore riprova!} \\
            }
        }
    }
    & & \\ \\

    \theutterance \stepcounter{utterance}  
    & & \multicolumn{4}{p{0.6\linewidth}}{
        \cellcolor[rgb]{0.9,0.9,0.9}{
            \makecell[{{p{\linewidth}}}]{
                \texttt{\tiny{[P1$\langle$GM]}}
                \texttt{Penalità: (4,1) non è vuoto, ma contiene 'D'.} \\
\\ 
\texttt{Avete accumulato collettivamente 2 di 16 penalità.} \\
\texttt{Per favore riprova!} \\
            }
        }
    }
    & \\ \\

    \theutterance \stepcounter{utterance}  
    & \multicolumn{4}{p{0.6\linewidth}}{
        \cellcolor[rgb]{0.9,0.9,0.9}{
            \makecell[{{p{\linewidth}}}]{
                \texttt{\tiny{[P1$\rangle$GM]}}
                \texttt{MUOVERE: D, (4,1)} \\
            }
        }
    }
    & & \\ \\

    \theutterance \stepcounter{utterance}  
    & & & \multicolumn{2}{p{0.3\linewidth}}{
        \cellcolor[rgb]{0.9,0.9,0.9}{
            \makecell[{{p{\linewidth}}}]{
                \texttt{\tiny{[GM$|$GM]}}
                \texttt{MUOVERE: D, (4,1)} \\
            }
        }
    }
    & & \\ \\

    \theutterance \stepcounter{utterance}  
    & & & \multicolumn{2}{p{0.3\linewidth}}{
        \cellcolor[rgb]{0.9,0.9,0.9}{
            \makecell[{{p{\linewidth}}}]{
                \texttt{\tiny{[GM$|$GM]}}
                \texttt{Penalità: (4,1) non è vuoto, ma contiene 'D'.} \\
\\ 
\texttt{Avete accumulato collettivamente 3 di 16 penalità.} \\
\texttt{Per favore riprova!} \\
            }
        }
    }
    & & \\ \\

    \theutterance \stepcounter{utterance}  
    & & \multicolumn{4}{p{0.6\linewidth}}{
        \cellcolor[rgb]{0.9,0.9,0.9}{
            \makecell[{{p{\linewidth}}}]{
                \texttt{\tiny{[P1$\langle$GM]}}
                \texttt{Penalità: (4,1) non è vuoto, ma contiene 'D'.} \\
\\ 
\texttt{Avete accumulato collettivamente 3 di 16 penalità.} \\
\texttt{Per favore riprova!} \\
            }
        }
    }
    & \\ \\

    \theutterance \stepcounter{utterance}  
    & \multicolumn{4}{p{0.6\linewidth}}{
        \cellcolor[rgb]{0.9,0.9,0.9}{
            \makecell[{{p{\linewidth}}}]{
                \texttt{\tiny{[P1$\rangle$GM]}}
                \texttt{DIRE: Attenzione: nel mio tabellone D è già a (3,1) e T è già a (4,1). Vuoi cambiare obiettivo per D e T, oppure confermi queste posizioni?} \\
            }
        }
    }
    & & \\ \\

    \theutterance \stepcounter{utterance}  
    & & & \multicolumn{2}{p{0.3\linewidth}}{
        \cellcolor[rgb]{0.9,0.9,0.9}{
            \makecell[{{p{\linewidth}}}]{
                \texttt{\tiny{[GM$|$GM]}}
                \texttt{DIRE: Attenzione: nel mio tabellone D è già a (3,1) e T è già a (4,1). Vuoi cambiare obiettivo per D e T, oppure confermi queste posizioni?} \\
            }
        }
    }
    & & \\ \\

    \theutterance \stepcounter{utterance}  
    & & \multicolumn{4}{p{0.6\linewidth}}{
        \cellcolor[rgb]{0.9,0.9,0.9}{
            \makecell[{{p{\linewidth}}}]{
                \texttt{\tiny{[GM$\rangle$P2]}}
                \texttt{Il tuo messaggio è stato trasmesso all'altro giocatore.} \\
\\ 
\\ 
\texttt{Stai giocando il round 2 di un massimo di 28.} \\
\texttt{Avete accumulato collettivamente 3 di 16 penalità.} \\
\texttt{L'altro giocatore ha inviato questo messaggio:} \\
\\ 
\\ 
\texttt{"Attenzione: nel mio tabellone D è già a (3,1) e T è già a (4,1). Vuoi cambiare obiettivo per D e T, oppure confermi queste posizioni?"} \\
\\ 
\\ 
\texttt{Qual è il tuo prossimo comando?} \\
            }
        }
    }
    & \\ \\

    \theutterance \stepcounter{utterance}  
    & & & \multicolumn{4}{p{0.6\linewidth}}{
        \cellcolor[rgb]{0.9,0.9,0.9}{
            \makecell[{{p{\linewidth}}}]{
                \texttt{\tiny{[GM$\langle$P2]}}
                \texttt{DIRE: Confermo D=(3,1) e T=(4,1). Nel mio tabellone sono vuote, quindi va bene. Procedo con lo spostamento successivo.} \\
            }
        }
    }
     \\ \\

    \theutterance \stepcounter{utterance}  
    & & & \multicolumn{2}{p{0.3\linewidth}}{
        \cellcolor[rgb]{0.9,0.9,0.9}{
            \makecell[{{p{\linewidth}}}]{
                \texttt{\tiny{[GM$|$GM]}}
                \texttt{DIRE: Confermo D=(3,1) e T=(4,1). Nel mio tabellone sono vuote, quindi va bene. Procedo con lo spostamento successivo.} \\
            }
        }
    }
    & & \\ \\

    \theutterance \stepcounter{utterance}  
    & & \multicolumn{4}{p{0.6\linewidth}}{
        \cellcolor[rgb]{0.9,0.9,0.9}{
            \makecell[{{p{\linewidth}}}]{
                \texttt{\tiny{[P1$\langle$GM]}}
                \texttt{Il tuo messaggio è stato trasmesso all'altro giocatore.} \\
\\ 
\\ 
\texttt{Stai giocando il round 2 di un massimo di 28.} \\
\texttt{Avete accumulato collettivamente 3 di 16 penalità.} \\
\texttt{L'altro giocatore ha inviato questo messaggio:} \\
\\ 
\\ 
\texttt{"Confermo D=(3,1) e T=(4,1). Nel mio tabellone sono vuote, quindi va bene. Procedo con lo spostamento successivo."} \\
\\ 
\\ 
\texttt{Qual è il tuo prossimo comando?} \\
            }
        }
    }
    & \\ \\

    \theutterance \stepcounter{utterance}  
    & \multicolumn{4}{p{0.6\linewidth}}{
        \cellcolor[rgb]{0.9,0.9,0.9}{
            \makecell[{{p{\linewidth}}}]{
                \texttt{\tiny{[P1$\rangle$GM]}}
                \texttt{MUOVERE: A, (7,7)} \\
            }
        }
    }
    & & \\ \\

    \theutterance \stepcounter{utterance}  
    & & & \multicolumn{2}{p{0.3\linewidth}}{
        \cellcolor[rgb]{0.9,0.9,0.9}{
            \makecell[{{p{\linewidth}}}]{
                \texttt{\tiny{[GM$|$GM]}}
                \texttt{MUOVERE: A, (7,7)} \\
            }
        }
    }
    & & \\ \\

    \theutterance \stepcounter{utterance}  
    & & & \multicolumn{2}{p{0.3\linewidth}}{
        \cellcolor[rgb]{0.9,0.9,0.9}{
            \makecell[{{p{\linewidth}}}]{
                \texttt{\tiny{[GM$|$GM]}}
                \texttt{Penalità: (7,7) non è vuoto, ma contiene 'A'.} \\
\\ 
\texttt{Avete accumulato collettivamente 4 di 16 penalità.} \\
\texttt{Per favore riprova!} \\
            }
        }
    }
    & & \\ \\

    \theutterance \stepcounter{utterance}  
    & & \multicolumn{4}{p{0.6\linewidth}}{
        \cellcolor[rgb]{0.9,0.9,0.9}{
            \makecell[{{p{\linewidth}}}]{
                \texttt{\tiny{[P1$\langle$GM]}}
                \texttt{Penalità: (7,7) non è vuoto, ma contiene 'A'.} \\
\\ 
\texttt{Avete accumulato collettivamente 4 di 16 penalità.} \\
\texttt{Per favore riprova!} \\
            }
        }
    }
    & \\ \\

    \theutterance \stepcounter{utterance}  
    & \multicolumn{4}{p{0.6\linewidth}}{
        \cellcolor[rgb]{0.9,0.9,0.9}{
            \makecell[{{p{\linewidth}}}]{
                \texttt{\tiny{[P1$\rangle$GM]}}
                \texttt{MUOVERE: M, (1,5)} \\
            }
        }
    }
    & & \\ \\

    \theutterance \stepcounter{utterance}  
    & & & \multicolumn{2}{p{0.3\linewidth}}{
        \cellcolor[rgb]{0.9,0.9,0.9}{
            \makecell[{{p{\linewidth}}}]{
                \texttt{\tiny{[GM$|$GM]}}
                \texttt{MUOVERE: M, (1,5)} \\
            }
        }
    }
    & & \\ \\

    \theutterance \stepcounter{utterance}  
    & & & \multicolumn{2}{p{0.3\linewidth}}{
        \cellcolor[rgb]{0.9,0.9,0.9}{
            \makecell[{{p{\linewidth}}}]{
                \texttt{\tiny{[GM$|$GM]}}
                \texttt{Penalità: (1,5) non è vuoto, ma contiene 'M'.} \\
\\ 
\texttt{Avete accumulato collettivamente 5 di 16 penalità.} \\
\texttt{Per favore riprova!} \\
            }
        }
    }
    & & \\ \\

    \theutterance \stepcounter{utterance}  
    & & \multicolumn{4}{p{0.6\linewidth}}{
        \cellcolor[rgb]{0.9,0.9,0.9}{
            \makecell[{{p{\linewidth}}}]{
                \texttt{\tiny{[P1$\langle$GM]}}
                \texttt{Penalità: (1,5) non è vuoto, ma contiene 'M'.} \\
\\ 
\texttt{Avete accumulato collettivamente 5 di 16 penalità.} \\
\texttt{Per favore riprova!} \\
            }
        }
    }
    & \\ \\

    \theutterance \stepcounter{utterance}  
    & \multicolumn{4}{p{0.6\linewidth}}{
        \cellcolor[rgb]{0.9,0.9,0.9}{
            \makecell[{{p{\linewidth}}}]{
                \texttt{\tiny{[P1$\rangle$GM]}}
                \texttt{DIRE: Scusa, sto sbagliando: A e M sono già in posizione. Faccio uno riepilogo delle mie posizioni attuali: P=(4,5), S=(1,7), D=(3,1), T=(4,1), A=(7,7), O=(7,4), M=(1,5). Risulta che tutti i miei oggetti sono già nelle posizioni obiettivo! Sei d'accordo?} \\
            }
        }
    }
    & & \\ \\

    \theutterance \stepcounter{utterance}  
    & & & \multicolumn{2}{p{0.3\linewidth}}{
        \cellcolor[rgb]{0.9,0.9,0.9}{
            \makecell[{{p{\linewidth}}}]{
                \texttt{\tiny{[GM$|$GM]}}
                \texttt{DIRE: Scusa, sto sbagliando: A e M sono già in posizione. Faccio uno riepilogo delle mie posizioni attuali: P=(4,5), S=(1,7), D=(3,1), T=(4,1), A=(7,7), O=(7,4), M=(1,5). Risulta che tutti i miei oggetti sono già nelle posizioni obiettivo! Sei d'accordo?} \\
            }
        }
    }
    & & \\ \\

    \theutterance \stepcounter{utterance}  
    & & \multicolumn{4}{p{0.6\linewidth}}{
        \cellcolor[rgb]{0.9,0.9,0.9}{
            \makecell[{{p{\linewidth}}}]{
                \texttt{\tiny{[GM$\rangle$P2]}}
                \texttt{Il tuo messaggio è stato trasmesso all'altro giocatore.} \\
\\ 
\\ 
\texttt{Stai giocando il round 3 di un massimo di 28.} \\
\texttt{Avete accumulato collettivamente 5 di 16 penalità.} \\
\texttt{L'altro giocatore ha inviato questo messaggio:} \\
\\ 
\\ 
\texttt{"Scusa, sto sbagliando: A e M sono già in posizione. Faccio uno riepilogo delle mie posizioni attuali: P=(4,5), S=(1,7), D=(3,1), T=(4,1), A=(7,7), O=(7,4), M=(1,5). Risulta che tutti i miei oggetti sono già nelle posizioni obiettivo! Sei d'accordo?"} \\
\\ 
\\ 
\texttt{Qual è il tuo prossimo comando?} \\
            }
        }
    }
    & \\ \\

    \theutterance \stepcounter{utterance}  
    & & & \multicolumn{4}{p{0.6\linewidth}}{
        \cellcolor[rgb]{0.9,0.9,0.9}{
            \makecell[{{p{\linewidth}}}]{
                \texttt{\tiny{[GM$\langle$P2]}}
                \texttt{DIRE: Capisco, tu sei già a posto. Ora devo allineare i miei oggetti. M, P, S, A e O sono già nelle posizioni obiettivo. Devo spostare solo D da (5,7) a (3,1) e T da (3,1) a (4,1). Inizio con T: (3,1) {-}> (4,1).} \\
            }
        }
    }
     \\ \\

    \theutterance \stepcounter{utterance}  
    & & & \multicolumn{2}{p{0.3\linewidth}}{
        \cellcolor[rgb]{0.9,0.9,0.9}{
            \makecell[{{p{\linewidth}}}]{
                \texttt{\tiny{[GM$|$GM]}}
                \texttt{DIRE: Capisco, tu sei già a posto. Ora devo allineare i miei oggetti. M, P, S, A e O sono già nelle posizioni obiettivo. Devo spostare solo D da (5,7) a (3,1) e T da (3,1) a (4,1). Inizio con T: (3,1) {-}> (4,1).} \\
            }
        }
    }
    & & \\ \\

    \theutterance \stepcounter{utterance}  
    & & \multicolumn{4}{p{0.6\linewidth}}{
        \cellcolor[rgb]{0.9,0.9,0.9}{
            \makecell[{{p{\linewidth}}}]{
                \texttt{\tiny{[P1$\langle$GM]}}
                \texttt{Il tuo messaggio è stato trasmesso all'altro giocatore.} \\
\\ 
\\ 
\texttt{Stai giocando il round 3 di un massimo di 28.} \\
\texttt{Avete accumulato collettivamente 5 di 16 penalità.} \\
\texttt{L'altro giocatore ha inviato questo messaggio:} \\
\\ 
\\ 
\texttt{"Capisco, tu sei già a posto. Ora devo allineare i miei oggetti. M, P, S, A e O sono già nelle posizioni obiettivo. Devo spostare solo D da (5,7) a (3,1) e T da (3,1) a (4,1). Inizio con T: (3,1) {-}> (4,1)."} \\
\\ 
\\ 
\texttt{Qual è il tuo prossimo comando?} \\
            }
        }
    }
    & \\ \\

    \theutterance \stepcounter{utterance}  
    & \multicolumn{4}{p{0.6\linewidth}}{
        \cellcolor[rgb]{0.9,0.9,0.9}{
            \makecell[{{p{\linewidth}}}]{
                \texttt{\tiny{[P1$\rangle$GM]}}
                \texttt{DIRE: Perfetto, io ho finito. Aspetto che tu completi i tuoi spostamenti. Fammi sapere quando sei pronto!} \\
            }
        }
    }
    & & \\ \\

    \theutterance \stepcounter{utterance}  
    & & & \multicolumn{2}{p{0.3\linewidth}}{
        \cellcolor[rgb]{0.9,0.9,0.9}{
            \makecell[{{p{\linewidth}}}]{
                \texttt{\tiny{[GM$|$GM]}}
                \texttt{DIRE: Perfetto, io ho finito. Aspetto che tu completi i tuoi spostamenti. Fammi sapere quando sei pronto!} \\
            }
        }
    }
    & & \\ \\

    \theutterance \stepcounter{utterance}  
    & & \multicolumn{4}{p{0.6\linewidth}}{
        \cellcolor[rgb]{0.9,0.9,0.9}{
            \makecell[{{p{\linewidth}}}]{
                \texttt{\tiny{[GM$\rangle$P2]}}
                \texttt{Il tuo messaggio è stato trasmesso all'altro giocatore.} \\
\\ 
\\ 
\texttt{Stai giocando il round 4 di un massimo di 28.} \\
\texttt{Avete accumulato collettivamente 5 di 16 penalità.} \\
\texttt{L'altro giocatore ha inviato questo messaggio:} \\
\\ 
\\ 
\texttt{"Perfetto, io ho finito. Aspetto che tu completi i tuoi spostamenti. Fammi sapere quando sei pronto!"} \\
\\ 
\\ 
\texttt{Qual è il tuo prossimo comando?} \\
            }
        }
    }
    & \\ \\

    \theutterance \stepcounter{utterance}  
    & & & \multicolumn{4}{p{0.6\linewidth}}{
        \cellcolor[rgb]{0.9,0.9,0.9}{
            \makecell[{{p{\linewidth}}}]{
                \texttt{\tiny{[GM$\langle$P2]}}
                \texttt{MUOVERE: T, (4, 1)} \\
            }
        }
    }
     \\ \\

    \theutterance \stepcounter{utterance}  
    & & & \multicolumn{2}{p{0.3\linewidth}}{
        \cellcolor[rgb]{0.9,0.9,0.9}{
            \makecell[{{p{\linewidth}}}]{
                \texttt{\tiny{[GM$|$GM]}}
                \texttt{MUOVERE: T, (4, 1)} \\
            }
        }
    }
    & & \\ \\

    \theutterance \stepcounter{utterance}  
    & & & \multicolumn{2}{p{0.3\linewidth}}{
        \cellcolor[rgb]{0.9,0.9,0.9}{
            \makecell[{{p{\linewidth}}}]{
                \texttt{\tiny{[GM$|$GM]}}
                \texttt{Penalità: (4,1) non è vuoto, ma contiene 'T'.} \\
\\ 
\texttt{Avete accumulato collettivamente 6 di 16 penalità.} \\
\texttt{Per favore riprova!} \\
            }
        }
    }
    & & \\ \\

    \theutterance \stepcounter{utterance}  
    & & \multicolumn{4}{p{0.6\linewidth}}{
        \cellcolor[rgb]{0.9,0.9,0.9}{
            \makecell[{{p{\linewidth}}}]{
                \texttt{\tiny{[GM$\rangle$P2]}}
                \texttt{Penalità: (4,1) non è vuoto, ma contiene 'T'.} \\
\\ 
\texttt{Avete accumulato collettivamente 6 di 16 penalità.} \\
\texttt{Per favore riprova!} \\
            }
        }
    }
    & \\ \\

    \theutterance \stepcounter{utterance}  
    & & & \multicolumn{4}{p{0.6\linewidth}}{
        \cellcolor[rgb]{0.9,0.9,0.9}{
            \makecell[{{p{\linewidth}}}]{
                \texttt{\tiny{[GM$\langle$P2]}}
                \texttt{DIRE: Scusa, nel mio tabellone T è già a (3,1) e non a (4,1). Devo solo spostare D da (5,7) a (3,1). Procedo ora con D.} \\
            }
        }
    }
     \\ \\

    \theutterance \stepcounter{utterance}  
    & & & \multicolumn{2}{p{0.3\linewidth}}{
        \cellcolor[rgb]{0.9,0.9,0.9}{
            \makecell[{{p{\linewidth}}}]{
                \texttt{\tiny{[GM$|$GM]}}
                \texttt{DIRE: Scusa, nel mio tabellone T è già a (3,1) e non a (4,1). Devo solo spostare D da (5,7) a (3,1). Procedo ora con D.} \\
            }
        }
    }
    & & \\ \\

    \theutterance \stepcounter{utterance}  
    & & \multicolumn{4}{p{0.6\linewidth}}{
        \cellcolor[rgb]{0.9,0.9,0.9}{
            \makecell[{{p{\linewidth}}}]{
                \texttt{\tiny{[P1$\langle$GM]}}
                \texttt{Il tuo messaggio è stato trasmesso all'altro giocatore.} \\
\\ 
\\ 
\texttt{Stai giocando il round 4 di un massimo di 28.} \\
\texttt{Avete accumulato collettivamente 6 di 16 penalità.} \\
\texttt{L'altro giocatore ha inviato questo messaggio:} \\
\\ 
\\ 
\texttt{"Scusa, nel mio tabellone T è già a (3,1) e non a (4,1). Devo solo spostare D da (5,7) a (3,1). Procedo ora con D."} \\
\\ 
\\ 
\texttt{Qual è il tuo prossimo comando?} \\
            }
        }
    }
    & \\ \\

    \theutterance \stepcounter{utterance}  
    & \multicolumn{4}{p{0.6\linewidth}}{
        \cellcolor[rgb]{0.9,0.9,0.9}{
            \makecell[{{p{\linewidth}}}]{
                \texttt{\tiny{[P1$\rangle$GM]}}
                \texttt{DIRE: OK, aspetto che tu sposti D. Fammi sapere quando hai finito!} \\
            }
        }
    }
    & & \\ \\

    \theutterance \stepcounter{utterance}  
    & & & \multicolumn{2}{p{0.3\linewidth}}{
        \cellcolor[rgb]{0.9,0.9,0.9}{
            \makecell[{{p{\linewidth}}}]{
                \texttt{\tiny{[GM$|$GM]}}
                \texttt{DIRE: OK, aspetto che tu sposti D. Fammi sapere quando hai finito!} \\
            }
        }
    }
    & & \\ \\

    \theutterance \stepcounter{utterance}  
    & & \multicolumn{4}{p{0.6\linewidth}}{
        \cellcolor[rgb]{0.9,0.9,0.9}{
            \makecell[{{p{\linewidth}}}]{
                \texttt{\tiny{[GM$\rangle$P2]}}
                \texttt{Il tuo messaggio è stato trasmesso all'altro giocatore.} \\
\\ 
\\ 
\texttt{Stai giocando il round 5 di un massimo di 28.} \\
\texttt{Avete accumulato collettivamente 6 di 16 penalità.} \\
\texttt{L'altro giocatore ha inviato questo messaggio:} \\
\\ 
\\ 
\texttt{"OK, aspetto che tu sposti D. Fammi sapere quando hai finito!"} \\
\\ 
\\ 
\texttt{Qual è il tuo prossimo comando?} \\
            }
        }
    }
    & \\ \\

    \theutterance \stepcounter{utterance}  
    & & & \multicolumn{4}{p{0.6\linewidth}}{
        \cellcolor[rgb]{0.9,0.9,0.9}{
            \makecell[{{p{\linewidth}}}]{
                \texttt{\tiny{[GM$\langle$P2]}}
                \texttt{MUOVERE: D, (3, 1)} \\
            }
        }
    }
     \\ \\

    \theutterance \stepcounter{utterance}  
    & & & \multicolumn{2}{p{0.3\linewidth}}{
        \cellcolor[rgb]{0.9,0.9,0.9}{
            \makecell[{{p{\linewidth}}}]{
                \texttt{\tiny{[GM$|$GM]}}
                \texttt{MUOVERE: D, (3, 1)} \\
            }
        }
    }
    & & \\ \\

    \theutterance \stepcounter{utterance}  
    & & & \multicolumn{2}{p{0.3\linewidth}}{
        \cellcolor[rgb]{0.9,0.9,0.9}{
            \makecell[{{p{\linewidth}}}]{
                \texttt{\tiny{[GM$|$GM]}}
                \texttt{D moved to (3, 1)} \\
            }
        }
    }
    & & \\ \\

    \theutterance \stepcounter{utterance}  
    & & \multicolumn{4}{p{0.6\linewidth}}{
        \cellcolor[rgb]{0.9,0.9,0.9}{
            \makecell[{{p{\linewidth}}}]{
                \texttt{\tiny{[P1$\langle$GM]}}
                \texttt{Il tuo messaggio è stato trasmesso all'altro giocatore.} \\
\\ 
\\ 
\texttt{Stai giocando il round 5 di un massimo di 28.} \\
\texttt{Avete accumulato collettivamente 6 di 16 penalità.} \\
\texttt{L'altro giocatore ha spostato un oggetto sulla sua griglia.} \\
\\ 
\\ 
\texttt{Qual è il tuo prossimo comando?} \\
            }
        }
    }
    & \\ \\

    \theutterance \stepcounter{utterance}  
    & \multicolumn{4}{p{0.6\linewidth}}{
        \cellcolor[rgb]{0.9,0.9,0.9}{
            \makecell[{{p{\linewidth}}}]{
                \texttt{\tiny{[P1$\rangle$GM]}}
                \texttt{DIRE: finito?} \\
            }
        }
    }
    & & \\ \\

    \theutterance \stepcounter{utterance}  
    & & & \multicolumn{2}{p{0.3\linewidth}}{
        \cellcolor[rgb]{0.9,0.9,0.9}{
            \makecell[{{p{\linewidth}}}]{
                \texttt{\tiny{[GM$|$GM]}}
                \texttt{DIRE: finito?} \\
            }
        }
    }
    & & \\ \\

    \theutterance \stepcounter{utterance}  
    & & \multicolumn{4}{p{0.6\linewidth}}{
        \cellcolor[rgb]{0.9,0.9,0.9}{
            \makecell[{{p{\linewidth}}}]{
                \texttt{\tiny{[GM$\rangle$P2]}}
                \texttt{Spostato 'D' in (3,1) con successo. La tua griglia aggiornata appare così:} \\
\texttt{```} \\
\texttt{ 1234567} \\
\texttt{╔═╤═══╤═╗} \\
\texttt{║◌│DT◌│◌║ 1} \\
\texttt{╟─┼──┬┘◌║ 2} \\
\texttt{║A│M◌│◌◌║ 3} \\
\texttt{╟─┴──┘◌S║ 4} \\
\texttt{║◌◌◌◌◌◌◌║ 5} \\
\texttt{║◌P◌◌◌◌◌║ 6} \\
\texttt{║◌◌O◌◌◌◌║ 7} \\
\texttt{╚═══════╝} \\
\texttt{```} \\
\texttt{Stai giocando il round 6 di un massimo di 28.} \\
\texttt{Avete accumulato collettivamente 6 di 16 penalità.} \\
\texttt{L'altro giocatore ha inviato questo messaggio:} \\
\\ 
\\ 
\texttt{"finito?"} \\
\\ 
\\ 
\texttt{Qual è il tuo prossimo comando?} \\
            }
        }
    }
    & \\ \\

    \theutterance \stepcounter{utterance}  
    & & & \multicolumn{4}{p{0.6\linewidth}}{
        \cellcolor[rgb]{0.9,0.9,0.9}{
            \makecell[{{p{\linewidth}}}]{
                \texttt{\tiny{[GM$\langle$P2]}}
                \texttt{DIRE: finito!} \\
            }
        }
    }
     \\ \\

    \theutterance \stepcounter{utterance}  
    & & & \multicolumn{2}{p{0.3\linewidth}}{
        \cellcolor[rgb]{0.9,0.9,0.9}{
            \makecell[{{p{\linewidth}}}]{
                \texttt{\tiny{[GM$|$GM]}}
                \texttt{DIRE: finito!} \\
            }
        }
    }
    & & \\ \\

    \theutterance \stepcounter{utterance}  
    & & & \multicolumn{2}{p{0.3\linewidth}}{
        \cellcolor[rgb]{0.9,0.9,0.9}{
            \makecell[{{p{\linewidth}}}]{
                \texttt{\tiny{[GM$|$GM]}}
                \texttt{true} \\
            }
        }
    }
    & & \\ \\

    \theutterance \stepcounter{utterance}  
    & & & \multicolumn{2}{p{0.3\linewidth}}{
        \cellcolor[rgb]{0.9,0.9,0.9}{
            \makecell[{{p{\linewidth}}}]{
                \texttt{\tiny{[GM$|$GM]}}
                \texttt{Initial states:} \\
\texttt{```} \\
\texttt{Player 1:} \\
\texttt{ 1234567} \\
\texttt{╔═╤═══╤═╗} \\
\texttt{║◌│◌DT│◌║ 1} \\
\texttt{╟─┼──┬┘◌║ 2} \\
\texttt{║◌│◌◌│◌◌║ 3} \\
\texttt{╟─┴──┘◌O║ 4} \\
\texttt{║M◌◌P◌◌◌║ 5} \\
\texttt{║◌◌◌◌◌◌◌║ 6} \\
\texttt{║S◌◌◌◌◌A║ 7} \\
\texttt{╚═══════╝} \\
\\ 
\\ 
\texttt{Player 2:} \\
\texttt{ 1234567} \\
\texttt{╔═╤═══╤═╗} \\
\texttt{║◌│◌T◌│◌║ 1} \\
\texttt{╟─┼──┬┘◌║ 2} \\
\texttt{║A│M◌│◌◌║ 3} \\
\texttt{╟─┴──┘◌S║ 4} \\
\texttt{║◌◌◌◌◌◌◌║ 5} \\
\texttt{║◌P◌◌◌◌◌║ 6} \\
\texttt{║◌◌O◌D◌◌║ 7} \\
\texttt{╚═══════╝} \\
\\ 
\texttt{```} \\
            }
        }
    }
    & & \\ \\

    \theutterance \stepcounter{utterance}  
    & & & \multicolumn{2}{p{0.3\linewidth}}{
        \cellcolor[rgb]{0.9,0.9,0.9}{
            \makecell[{{p{\linewidth}}}]{
                \texttt{\tiny{[GM$|$GM]}}
                \texttt{End states:} \\
\texttt{```} \\
\texttt{Player 1:} \\
\texttt{ 1234567} \\
\texttt{╔═╤═══╤═╗} \\
\texttt{║◌│◌DT│◌║ 1} \\
\texttt{╟─┼──┬┘◌║ 2} \\
\texttt{║◌│◌◌│◌◌║ 3} \\
\texttt{╟─┴──┘◌O║ 4} \\
\texttt{║M◌◌P◌◌◌║ 5} \\
\texttt{║◌◌◌◌◌◌◌║ 6} \\
\texttt{║S◌◌◌◌◌A║ 7} \\
\texttt{╚═══════╝} \\
\\ 
\\ 
\texttt{Player 2:} \\
\texttt{ 1234567} \\
\texttt{╔═╤═══╤═╗} \\
\texttt{║◌│DT◌│◌║ 1} \\
\texttt{╟─┼──┬┘◌║ 2} \\
\texttt{║A│M◌│◌◌║ 3} \\
\texttt{╟─┴──┘◌S║ 4} \\
\texttt{║◌◌◌◌◌◌◌║ 5} \\
\texttt{║◌P◌◌◌◌◌║ 6} \\
\texttt{║◌◌O◌◌◌◌║ 7} \\
\texttt{╚═══════╝} \\
\\ 
\texttt{```} \\
            }
        }
    }
    & & \\ \\

    \theutterance \stepcounter{utterance}  
    & & & \multicolumn{2}{p{0.3\linewidth}}{
        \cellcolor[rgb]{0.9,0.9,0.9}{
            \makecell[{{p{\linewidth}}}]{
                \texttt{\tiny{[GM$|$GM]}}
                \texttt{* success: True} \\
\texttt{* lose: False} \\
\texttt{* aborted: False} \\
\texttt{{-}{-}{-}{-}{-}{-}{-}} \\
\texttt{* Shifts: 0.00} \\
\texttt{* Max Shifts: 12.00} \\
\texttt{* Min Shifts: 6.00} \\
\texttt{* End Distance Sum: 25.98} \\
\texttt{* Init Distance Sum: 31.07} \\
\texttt{* Expected Distance Sum: 29.33} \\
\texttt{* Penalties: 6.00} \\
\texttt{* Max Penalties: 16.00} \\
\texttt{* Rounds: 6.00} \\
\texttt{* Max Rounds: 28.00} \\
\texttt{* Object Count: 7.00} \\
            }
        }
    }
    & & \\ \\

\end{supertabular}
}

\end{document}
